\documentclass[11pt]{article}

\usepackage{amsmath}
\usepackage{amssymb}
\usepackage{enumitem}
\usepackage{booktabs}
\usepackage{fancyhdr}
\usepackage[margin=1in]{geometry}
\usepackage{hyperref}
\usepackage{xcolor}

\hypersetup{
  colorlinks=true,
  linkcolor=blue,
  urlcolor=blue,
}

\pagestyle{fancy}
\fancyhf{}
\lhead{CEC 315: Signals and Systems}
\rhead{Spring 2026}
\cfoot{\thepage}

\newcommand{\boxedeq}[1]{\boxed{\,#1\,}}

\title{CEC 315 --- Lecture 23 Feedback Systems: Official Solutions}
\author{CEC 315 Signals and Systems \\ Spring 2026}
\date{}

\begin{document}
\maketitle

\noindent\textbf{Topic:} Lecture 23 --- Feedback systems, block diagram simplification,
closed-loop poles, stability, gain/phase margins, Nyquist. \\
\textbf{Exam coverage:} Exam 3. \\
\textbf{Source:} Transcribed from \texttt{lctr23-exercise-solutions.pdf} (instructor solutions).

\bigskip\hrule\bigskip

\section*{Block Diagram Simplification}

\subsection*{Problem 1 --- Cascade + Feedback}

The forward path is two blocks in cascade:
\[
H(s) = \frac{2}{s+1}\cdot\frac{1}{s+4} = \frac{2}{(s+1)(s+4)}.
\]
The feedback path is unity ($G=1$). Apply the feedback formula $Q = H/(1+GH)$:
\[
Q(s) = \frac{H(s)}{1 + G\cdot H(s)} = \frac{\dfrac{2}{(s+1)(s+4)}}{1 + \dfrac{2}{(s+1)(s+4)}}.
\]
Multiply numerator and denominator by $(s+1)(s+4)$:
\[
Q(s) = \frac{2}{(s+1)(s+4)+2} = \frac{2}{s^2+5s+4+2} = \frac{2}{s^2+5s+6} = \frac{2}{(s+2)(s+3)}.
\]
\[
\boxedeq{Q(s) = \frac{2}{s^2+5s+6} = \frac{2}{(s+2)(s+3)}}
\]
Closed-loop poles at $s=-2$ and $s=-3$. Both negative $\Longrightarrow$ \textbf{stable}.

\subsection*{Problem 2 --- Parallel + Feedback}

\textbf{Step 1 --- Parallel combination.} The two parallel blocks add:
\[
H(s) = \frac{3}{s+2} + \frac{1}{s+5}
     = \frac{3(s+5)+(s+2)}{(s+2)(s+5)}
     = \frac{4s+17}{(s+2)(s+5)}
     = \frac{4s+17}{s^2+7s+10}.
\]
\textbf{Step 2 --- Feedback.} Forward $=H(s)$, feedback $G=4$. Apply $Q = H/(1+GH)$:
\[
Q(s) = \frac{H(s)}{1+4H(s)}
     = \frac{\dfrac{4s+17}{s^2+7s+10}}{1+\dfrac{4(4s+17)}{s^2+7s+10}}
     = \frac{4s+17}{s^2+7s+10+16s+68}.
\]
\[
\boxedeq{Q(s) = \frac{4s+17}{s^2+23s+78}}
\]
\textbf{Check stability:} $s^2+23s+78=0$. Quadratic formula:
\[
s = \frac{-23\pm\sqrt{529-312}}{2} = \frac{-23\pm\sqrt{217}}{2} = \frac{-23\pm 14.73}{2}.
\]
So $s\approx -4.14$ and $s\approx -18.87$. Both negative $\Longrightarrow$ \textbf{stable}.

\subsection*{Problem 3 --- Nested Loops}

\textbf{Step 1 --- Inner loop.} Forward $=1/s$, feedback $=3$:
\[
Q_{\text{inner}}(s) = \frac{1/s}{1+3/s} = \frac{1/s}{(s+3)/s} = \frac{1}{s+3}.
\]
\textbf{Step 2 --- Outer loop.} Forward $=Q_{\text{inner}}=1/(s+3)$, feedback $=1$:
\[
Q(s) = \frac{1/(s+3)}{1+1/(s+3)} = \frac{1/(s+3)}{(s+4)/(s+3)} = \frac{1}{s+4}.
\]
\[
\boxedeq{Q(s) = \frac{1}{s+4}}
\]
Closed-loop pole at $s=-4$. \textbf{Stable.}

\textbf{Strategy reminder:} Always start with the innermost loop. Replace it with a single block.
Then treat the result as the forward path of the outer loop.

\bigskip\hrule\bigskip

\section*{Closed-Loop Poles and Stability}

\subsection*{Problem 4 --- Unity Feedback with $H(s)=K/(s+5)$}

\textbf{(a)} Forward: $H=K/(s+5)$, feedback: $G=1$. Then
\[
Q(s) = \frac{K/(s+5)}{1+K/(s+5)} = \frac{K}{s+5+K}.
\]
Closed-loop pole: $s = -(5+K)$.
\[
\boxedeq{Q(s) = \frac{K}{s+5+K}, \qquad s_{\text{cl}} = -(5+K)}
\]

\textbf{(b)} For CT stability (LHP): $5+K>0$, so
\[
\boxedeq{K>-5}.
\]

\textbf{(c)} We want $s=-10$: $-(5+K)=-10 \Longrightarrow \boxedeq{K=5}$.

\subsection*{Problem 5 --- Second-Order Feedback with $H(s)=K/[(s+1)(s+2)]$}

\textbf{(a)} Apply $Q=H/(1+GH)$ with $G=1$:
\[
Q(s) = \frac{K/[(s+1)(s+2)]}{1+K/[(s+1)(s+2)]} = \frac{K}{(s+1)(s+2)+K} = \frac{K}{s^2+3s+(2+K)}.
\]
Characteristic equation:
\[
\boxedeq{s^2 + 3s + (2+K) = 0}.
\]

\textbf{(b)} For a second-order polynomial $s^2+bs+c$ to have both roots in the LHP, we need $b>0$ and $c>0$.
\begin{itemize}[nosep]
  \item Here $b=3>0$ always.
  \item $c=2+K>0$ requires $K>-2$.
\end{itemize}
\[
\boxedeq{K>-2}.
\]

\textbf{Verification at the boundary:} At $K=-2$, the characteristic equation is
$s^2+3s = s(s+3)=0$, giving poles at $s=0$ and $s=-3$.
The pole at $s=0$ is on the imaginary axis, confirming marginal stability at the boundary.

\subsection*{Problem 6 --- DT Feedback with $H(z)=Kz^{-1}/(1-0.8z^{-1})$}

Forward: $H(z)=Kz^{-1}/(1-0.8z^{-1})$. Feedback: $G=1$.
\[
Q(z) = \frac{Kz^{-1}/(1-0.8z^{-1})}{1+Kz^{-1}/(1-0.8z^{-1})}
     = \frac{Kz^{-1}}{1-0.8z^{-1}+Kz^{-1}}
     = \frac{Kz^{-1}}{1+(K-0.8)z^{-1}}.
\]
Rewriting in terms of $z$:
\[
Q(z) = \frac{K}{z+(K-0.8)}.
\]

\textbf{(a)} Closed-loop pole: $z = -(K-0.8) = 0.8-K$.
\[
\boxedeq{z_{\text{cl}} = 0.8-K}.
\]

\textbf{(b)} DT stability: $|z_{\text{cl}}|<1$.
\[
|0.8-K|<1 \;\Longleftrightarrow\; -1<0.8-K<1 \;\Longleftrightarrow\; -0.2<K<1.8.
\]
For $K>0$:
\[
\boxedeq{0<K<1.8}.
\]
\textbf{Check:} At $K=1$, pole at $z=-0.2$ (inside, stable). At $K=1.8$, pole at $z=-1$
(on unit circle, marginally stable). At $K=2$, pole at $z=-1.2$ (outside, unstable).

\bigskip\hrule\bigskip

\section*{Gain and Phase Margins}

\subsection*{Problem 7 --- Margins from Bode Data}

\textbf{Given data:}
\begin{center}
\begin{tabular}{lcc}
\toprule
 & $|GH|$ (dB) & $\angle GH$ \\
\midrule
$\omega_2=5$ rad/s (where $|GH|=0$ dB) & $0$ dB & $-150^\circ$ \\
$\omega_1=20$ rad/s (where $\angle GH=-180^\circ$) & $-8$ dB & $-180^\circ$ \\
\bottomrule
\end{tabular}
\end{center}

\textbf{(a) Phase margin.} Use the 0-dB crossing $\omega_2 = \omega_{gc} = 5$ rad/s.
By definition $\text{PM}=180^\circ+\angle GH(\omega_{gc})$:
\[
\text{PM} = 180^\circ + \angle GH(\omega_2) = 180^\circ + (-150^\circ) = 30^\circ.
\]
\[
\boxedeq{\text{PM} = 30^\circ}.
\]

\textbf{(b) Gain margin.} Use the $-180^\circ$ phase-crossing $\omega_1 = 20$ rad/s.
By definition $\text{GM} = 0\text{ dB}-|GH|_{\text{dB}}$ at that frequency:
\[
\text{GM} = 0\text{ dB} - |GH(\omega_1)|_{\text{dB}} = 0-(-8) = 8\text{ dB}.
\]
\[
\boxedeq{\text{GM} = 8\text{ dB}}.
\]

\textbf{(c) Stability.} Both margins are positive $\Longrightarrow$ \textbf{stable}.
\[
\boxedeq{\text{Stable}}.
\]

\textbf{(d) Maximum tolerable time delay.} A pure delay of $\tau$ seconds adds phase $-\omega\tau$ radians
(no gain change). The binding constraint is at the gain-crossover frequency $\omega_{gc} = \omega_2 = 5$ rad/s,
where the system can tolerate $30^\circ$ (the phase margin) of extra negative phase before instability:
\[
\omega_2\cdot\tau_{\max} = 30^\circ\times\frac{\pi}{180^\circ} = 0.524\text{ rad} = \text{PM(rad)}.
\]
\[
\tau_{\max} = \frac{\text{PM(rad)}}{\omega_{gc}} = \frac{0.524}{5} = 0.105\text{ s} \approx 105\text{ ms}.
\]
\[
\boxedeq{\tau_{\max} \approx 0.105\text{ s} = 105\text{ ms}}.
\]

\bigskip\hrule\bigskip

\section*{True/False}

\subsection*{Problem 8 --- ``Negative feedback always stabilizes a system.''}

\[
\boxedeq{\textbf{FALSE}}
\]

\textbf{Explanation:} Negative feedback \emph{can} stabilize an unstable open-loop system, but it can also
destabilize a stable open-loop system if the loop gain is too high. The classical audio-feedback example
(squealing microphone) and Example 23.5 in the lecture (which goes unstable for $K>12$) both show that
enough gain in the loop pushes closed-loop poles into the RHP. Whether feedback stabilizes depends on the
loop transfer function $GH$ and the gain $K$, not on the sign of the summer alone.

\subsection*{Problem 9 --- ``If the phase margin is $60^\circ$, the system is stable.''}

\[
\boxedeq{\textbf{TRUE}}
\]

\textbf{Explanation:} A positive phase margin means that at the gain-crossover frequency the loop phase
has not yet reached $-180^\circ$, so the system has not reached the instability condition $GH=-1$.
In standard exam problems we also implicitly assume the gain margin is positive, and a positive
(in fact a generous $60^\circ$) phase margin indicates stability. $60^\circ$ is a common
``well-damped'' design target.

\subsection*{Problem 10 --- ``A Nyquist plot is constructed by evaluating $G(j\omega)H(j\omega)$ and plotting it in the complex plane.''}

\[
\boxedeq{\textbf{TRUE}}
\]

\textbf{Explanation:} At each frequency $\omega$, you evaluate the complex number $GH(j\omega)$
(magnitude and phase) and plot it as a point in the complex plane. Connecting all such points
for $\omega$ running from $-\infty$ to $+\infty$ traces out the Nyquist contour. Instability is
then diagnosed from encirclements of the critical point ($-1$ for unit loop gain, or $-1/K$
when the gain $K$ is factored out of $GH$).

\bigskip\hrule\bigskip

\section*{Key Takeaways}

\begin{enumerate}[leftmargin=*]
  \item \textbf{Closed-loop formula:} $Q = \dfrac{H}{1+GH}$. The denominator $1+GH$ sets the
        closed-loop poles. For unity feedback ($G=1$), this reduces to $H/(1+H)$.
  \item \textbf{Block diagram simplification --- three rules:} cascade (multiply), parallel (add),
        feedback loop ($H/(1+GH)$). Always reduce innermost loops first.
  \item \textbf{Continuous-time stability:} all closed-loop poles must lie in the
        \emph{left-half plane} ($\operatorname{Re}\{s\}<0$).
  \item \textbf{Discrete-time stability:} all closed-loop poles must lie \emph{inside the unit circle}
        ($|z|<1$).
  \item \textbf{Characteristic equation:} set $1+GH=0$ (equivalently, closed-loop denominator $=0$);
        this is how the stable range of $K$ is found symbolically.
  \item \textbf{Instability condition:} $GH=-1$, i.e.\ $|GH|=0$ dB \emph{and}
        $\angle GH=-180^\circ$ simultaneously.
  \item \textbf{Phase margin (PM):} measured at the 0-dB gain-crossover frequency $\omega_{gc}$:
        \[\text{PM} = 180^\circ + \angle GH(\omega_{gc}).\]
  \item \textbf{Gain margin (GM):} measured at the $-180^\circ$ phase-crossover frequency $\omega_{pc}$:
        \[\text{GM} = 0\text{ dB} - |GH(\omega_{pc})|_{\text{dB}} = -|GH(\omega_{pc})|_{\text{dB}}.\]
  \item \textbf{Max tolerable time delay (from PM):}
        \[\tau_{\max} = \frac{\text{PM (in radians)}}{\omega_{gc}}.\]
  \item \textbf{Nyquist criterion:} plot $GH(j\omega)$ in the complex plane. For a stable open-loop system,
        closed-loop stability requires that the Nyquist curve not encircle the critical point $-1/K$.
        Encirclements of $-1/K$ indicate RHP closed-loop poles.
\end{enumerate}

\end{document}
